\subsubsection{The ramified column lattice of a selected word}
\label{sec:selected-ramified-donor}

At a fixed translate of the selected family take \(w\in W\), put
\(T=T_w\) and \(q=T-1\), and write \(M=M_w\) for the full
stage-zero class lattice. We work over the coefficient field \(K\)
of this translate. The selection identifies the point, its object
carrying the arithmetic parameters \(L\), and its bundle witness as
one dependent triple at a constructed point. The point and the
true-column condition are obtained from the same local lattice
record. The comparison with the geometric realization preserves
the selected word, its column coordinates, the level operators and
the period data. This choice is available at each index of the
family, including the translated indices.
\releasecite{Development/AbelianVarietyTheory/Definitions/GlobalTwisting/ValidForConstructedPointClassesEventuallyWorking.lean}{250}{Selection of the point, object and witness},
\releasecite{Development/AbelianVarietyTheory/Proofs/GlobalInterpolation/ExistsObjectsOfAdmissibleTwistDatum.lean}{114}{The common local record},
\releasecite{Development/AutomorphicArithmetic/ThetaCorrespondence/ThetaPairLifts/ThirdPlaceSeries/TwistedBundle.lean}{57}{The true-column bundle},
\releasecite{Development/AutomorphicArithmetic/ThetaCorrespondence/ThetaPairLifts/ThirdPlaceSeries/LinkROfIsLinked.lean}{50}{The level and period data} and
\releasecite{Development/AbelianVarietyTheory/Proofs/IwasawaDecomposition/ExistsSelectedRealization.lean}{189}{The word and coordinates in the selected realization}.

Let \(\kappa_w\) be evaluation of a rational decorated class
function on the chosen representatives in the formal stage-zero
space \((D.\mathrm{fm}\,w).\mathrm{lat}\,0\).
The class covariance gives its inverse, as in the evaluation of
\(M_w\) above. Evaluation preserves integrality because the
formal lattice uses the same polynomial coefficient lattice
at each representative. It intertwines each diamond: taking
a translated point back to its representative applies the same
coefficient transport. The norm is the orthogonal weighted sum
of the normalized class forms in either description, at every
coefficient embedding.
\releasecite{Development/AbelianVarietyTheory/Proofs/DecoratedLatticeProofs/ColStep0MemLat.lean}{48}{The integral class test},
\releasecite{Development/AbelianVarietyTheory/Proofs/DecoratedLatticeProofs/EmbFColFormMemMram.lean}{51}{The diamond action under evaluation} and
\releasecite{Development/AbelianVarietyTheory/Proofs/DecoratedLatticeProofs/NrmStepZeroEq.lean}{54}{The class forms at all embeddings}.

Let \(S_a\) be the sum of the \(q\) diamonds in direction \(a\)
for \(1\le a\le f\), and put
\[
 V_{\rm ram}^{\rm class}=\bigcap_{a=1}^f\ker(S_a:M_K\to M_K),
 \qquad H=H_w=M\cap V_{\rm ram}^{\rm class}.
\]
The map \(j_w=D.\mathrm{nuIn}\,w\) includes the rational
space of \(D.\mathrm{div}.\mathrm{Nu}\,w\) into the formal
stage-zero space with image its ramified subspace. It induces
the stage-zero metric on that source. Since \(\kappa_w\)
intertwines the sums, the required rational comparison is
\begin{equation}
\label{eq:selected-donor-isometry}
 \mathcal E_w=\kappa_w^{-1}j_w:
     (D.\mathrm{div}.\mathrm{Nu}\,w).V
           \xrightarrow{\ \sim\ } V_{\rm ram}^{\rm class},
 \qquad
 \mathcal E_w(\text{\rm integral ramified pullback})=H_w .
\end{equation}
The second identity is the stage-zero integrality test on this
subspace. The comparison is an isometry of the Hermitian forms
at every embedding. In particular conjugate embeddings use the
conjugate class forms.
\releasecite{Development/AbelianVarietyTheory/Proofs/SatakeParameters/RamifiedColumnsTrueColumnsEqConcrete.lean}{18}{The comparison on the true columns},
\releasecite{Formalism/LocalRingCardinality/DataDIVDefs.lean}{12}{The integral pullback} and
\releasecite{Formalism/LocalRingCardinality/RhoDetPrimeEqMuTildeMulChiPow.lean}{291}{The isometric inclusion and its image}.

\paragraph{Integral rank and the lower mean.}
\leavevmode\teamanchor{proof.ramified-residue-orbits}%
The diamond group \((\mathbf F_T^\times)^f\) acts freely on
classes by the argument in the
\hyperlink{proof.ramified-column-loss}{ramified column count}.
Translate the integral fibre \(L_o\) at one chosen representative
of each free orbit \(o\). The integral Weyl lattice and determinant
factor are preserved in the transported frame. The finite-order
level character acts by a unit of absolute value one at every
embedding. The translates have the same class weight and
disjoint supports. Thus the following module identification
gives the orthogonal class metric at every embedding:
\[
 M=\bigoplus_o(\mathcal O_K^q)^{\otimes f}\otimes L_o,\qquad
 A_q=\ker(\Sigma:\mathcal O_K^q\to\mathcal O_K),\qquad
 H=\bigoplus_o A_q^{\otimes f}\otimes L_o .
\]
For the last equality split \(\mathcal O_K^q=A_q\oplus\mathcal O_K e_q\),
where \(e_q\) is the last coordinate vector and \(\Sigma\) the
coordinate sum. In direction \(a\), \(S_a\) factors as summation
followed by the diagonal inclusion. The latter has a left inverse
given by taking a coordinate. The simultaneous kernel therefore
retains exactly the summands with \(A_q\) in every position. The
saturation and the kernel description survive every base change,
including at places dividing \(q\). Writing
\(r_\diamond=\sum_o\operatorname{rk}L_o\), we obtain
\begin{equation}
\label{eq:selected-ramified-rank}
 d_M=q^f r_\diamond,\qquad d_r=(q-1)^f r_\diamond,\qquad
 \theta_w=1-d_r/d_M=1-(1-1/q)^f .
\end{equation}
\releasecite{Formalism/LocalToGlobal/GlobalTheory/GL6/DiagonalSlopes/SlopeRamifiedInducedPrimeField.lean}{653}{The slice and its translates},
\releasecite{Formalism/LocalToGlobal/GlobalTheory/GL6/DiagonalSlopes/SlopeRamifiedInducedPrimeField.lean}{731}{Integral membership},
\releasecite{Formalism/LocalToGlobal/GlobalTheory/GL6/DiagonalSlopes/SlopeRamifiedInducedPrimeField.lean}{809}{Orthogonality} and
\releasecite{Formalism/LocalToGlobal/GlobalTheory/GL6/DiagonalSlopes/SlopeFormulaOfCharacterFinite.lean}{17}{The character metric}.

\leavevmode\teamanchor{proof.ramified-kernel-lower-mean}%
Give \(Q=M/H\) its quotient metric. The sum operators induce an
integral injection \(Q\to M^{\oplus f}\). Each \(S_a\) has norm
at most \(q\) as a sum of \(q\) isometries. Taking the least-norm
lift in \(M\otimes_{\mathcal O_K,\sigma}\mathbb C\) at each
embedding \(\sigma\) gives a squared operator bound of \(fq^2\)
for the injection. For the orthogonal direct sum,
\(\mu_{\max}(M^{\oplus f})=\mu_{\max}(M)\). Indeed, filter a
saturated sublattice by the successive coordinate kernels. Each
quotient injects integrally and contractively into the next
coordinate. Passing to the saturation of its image can only
raise the degree. Additivity of degrees with quotient metrics
gives the upper bound and using a single coordinate gives equality.
The determinant inequality for the injection of \(Q\) now gives
\(\mu(Q)\le\mu_{\max}(M)+\log(fq^2)\), since its image lies in the
saturation in \(M^{\oplus f}\) and its squared determinant norm
is at most \((fq^2)^{\operatorname{rk}Q}\).
Degree additivity for \(H\subseteq M\) yields
\begin{equation}
\label{eq:ramified-quotient-bound}
 \mu(H_w)\ge\mu(M_w)-\delta_{\rm ram}(w),\qquad
 \delta_{\rm ram}(w)=\frac{\theta_w}{1-\theta_w}
  \bigl(\mu_{\max}(M_w)-\mu(M_w)+\log(fq^2)\bigr).
\end{equation}
Any rational subspace of \(H_K\) has the same saturated lattice
and induced metric in \(H\) and in \(M\). Hence
\(\mu_{\max}(H)\le\mu_{\max}(M)\).

Here \(\delta_{\rm ram}(w)=o(1)\) uniformly. Indeed
\(\theta_w\le f/(T_w-1)\), \(T_w\ge b^3\), and the
\hyperlink{proof.reference-schur-slopes}{reference estimates}
bound the spread \(\mu_{\max}(M_w)-\mu(M_w)\) by
\(O(N_Ta_E L_a(1+m/k))\), with \(a_E=xb\) and
\(L_a=\log(a_E\operatorname{rd}(E))\).
All the remaining factors and \(\log T_w\) are polynomially
bounded in \(x\) at fixed precision. This is the mean estimate
used in the \hyperlink{proof.divided-column-slopes}{divided-source slopes}.

\paragraph{The exact mean and the projected lattice.}
\leavevmode\teamanchor{proof.ramified-integral-kernel}%
The integral description also gives an exact formula. In the
one-factor basis \(e_i-e_q\), \(1\le i<q\), the Gram matrix is
\(G=I+\mathbf1\mathbf1^{\mathsf T}\) of determinant \(q\).
In the \(f\)-fold tensor the determinant is
\(q^{f(q-1)^{f-1}}\). The finite determinants of the orbit
fibres and the class metrics contribute the same amount per
rank on both sides, since each orbit retains the same rank
fraction. Thus the squared-norm convention gives
\begin{equation}
\label{eq:exact-ramified-column-slope}
 \mu(H_w)=\mu(M_w)-\frac{f\log(T_w-1)}{T_w-2}.
\end{equation}
This is the mean clause of the fitted \(\mathrm{adm26}\) input.
For comparison let \(P\) be the image of \(M\) under the product
of the rational orthogonal projections
\(I-q^{-1}\mathbf1\mathbf1^{\mathsf T}\) in the \(f\) directions.
In one factor the projected basis \(e_i-q^{-1}\mathbf1\),
\(1\le i<q\), has Gram \(G^{-1}\). The inclusion of the
difference basis has matrix \(G\) in this basis. Its integral
Smith factors are \(q\) once and \(1\) repeated \(q-2\) times,
because the determinant is \(q\) and the cokernel is generated
by the common class of the projected basis vectors. Tensor
the Smith bases and localize. The exponent \(j\operatorname{ord}_v(q)\)
has multiplicity \(r_\diamond\binom fj(q-2)^{f-j}\) for
\(0\le j\le f\), with multiplicities added when exponents
coincide. In particular,
\[
 H\subseteq P\subseteq q^{-f}H,\qquad
 \operatorname{length}_{\mathcal O_v}(P_v/H_v)
      =\frac{f d_r}{q-1}\operatorname{ord}_v(q),\qquad
 \mu(P)=\mu(M)+\frac{f\log q}{q-1}.
\]
At every \(v\mid T_w\) the integer \(q\) is a unit, so
\(P_v=H_v\). The donor index there is the Hecke index below.
\releasecite{Formalism/LocalToGlobal/Columns/SliceIntegralDeltaNeBot.lean}{858}{The exact concrete mean} and
\releasecite{Formalism/SharpRestriction/Divisibility/LocalQuadraticParts.lean}{621}{The formula in the fitted normalization}.
\rhrange{Equation~\eqref{eq:exact-ramified-column-slope} gives the exact
squared-metric mean of the integral ramified lattice.}
{The \hyperlink{proof.ramified-kernel-lower-mean}{quotient lower mean}
\(\mu(H_w)\ge\mu(M_w)-o(1)\) and
\(\mu_{\max}(H_w)\le\mu_{\max}(M_w)\), with the actual rank,
donor and local indices.}

\paragraph{The undivided period source.}
\leavevmode\teamanchor{proof.selected-undivided-donor}%
Let \(e_{\nu,w}\) denote the map from the fixed rational
column coordinates to \((D.\mathrm{div}.\mathrm{Nu}\,w).V\).
If \(z\) has class function \(\phi_z\) in those coordinates,
evaluation on every representative gives
\[
 j_w(e_{\nu,w}z)=\kappa_w(\phi_z),
 \qquad \mathcal E_w(e_{\nu,w}z)=\phi_z .
\]
Injectivity of \(j_w\) determines \(e_{\nu,w}\) from this equation.
The all-embedding Gram identity for this map follows from the
stage-zero isometry \eqref{eq:selected-donor-isometry}. With these
coordinates the division construction chooses the undivided lattice
as \(e_{\nu,w}^{-1}(L.\mathrm{NuLat}\,w)\). Consequently the
class-function lattice in the period descent is
\begin{equation}
\label{eq:selected-undivided-donor}
 N_w=\mathcal E_w(L.\mathrm{NuLat}\,w)
      \subseteq H_w\subseteq M_w .
\end{equation}
The selected divided coordinate lattice has image
\(I_w(L.\mathrm{NuLat}\,w)\) under \(e_{\nu,w}\), with
\(I_w=\prod_v\mathfrak p_v^{-\lfloor e_vg_{wv}\rfloor}\).
Transport commutes with multiplication by this fractional ideal.
Applying its inverse recovers the undivided lattice, so the
comparison from the final divided realization uses the same choice.
\releasecite{Development/AutomorphicArithmetic/ThetaCorrespondence/ThetaPairLifts/ThirdPlaceSeries/ExistsRowsKCColPins.lean}{146}{The tied coordinates and period data},
\releasecite{Development/AbelianVarietyTheory/Proofs/IwasawaDecomposition/ChKerCCongr.lean}{112}{The column identity and all-place Gram identity},
\releasecite{Development/AutomorphicArithmetic/ThetaCorrespondence/ThetaPairLifts/ThirdPlaceSeries/Wire4F.lean}{83}{Application at the selected index},
\releasecite{Formalism/LocalToGlobal/Divisibility/RecBridgeMkF.lean}{57}{The undivided lattice in coordinates},
\releasecite{Formalism/LocalToGlobal/Divisibility/DivisibilityConstruction/ShiftOfColTie.lean}{53}{The fractional-ideal shift} and
\releasecite{Formalism/LocalToGlobal/GroupLattices/JointServices/ShiftIdealSgFl.lean}{57}{Its cancellation}.

For every used tuple \(\boldsymbol\mu\), let \(e_{\boldsymbol\mu}\)
be the realization of its placed top row coordinates, and let
\(\mathcal P_{w,\boldsymbol\mu}(\phi,z)\) be the \(K\)-valued
top class sum of Section~\ref{sec:periods} in those coordinates.
On \(n\) in the formal rational column space and a placed test
row \(z\), the pairing identity is
\begin{equation}
\label{eq:selected-donor-placed-period}
 \operatorname{perCol}(w,\boldsymbol\mu)
       (j_N n,e_{\boldsymbol\mu}z)
   =\iota_K\bigl(\mathcal P_{w,\boldsymbol\mu}(\mathcal E_w n,z)\bigr),
\end{equation}
where \(j_N\) is the induced complex inclusion. At the base
tuple this is evaluation of the terminal class sum. At a moving
tuple, transport the initial class values by the prescribed
tuple map and apply the normalized step entries on the moved
representatives. The class evaluation and the coefficient
transport intertwine that map, so the sum gives the placed period
with its tuple normalization. At the base let \(e_R\) be the
identification of the descended row coordinates with the formal
ones. The period kills the projection kernel and the descended
rows span. Thus \eqref{eq:selected-donor-placed-period} and
injectivity of \(\iota_K\) give
\begin{equation}
\label{eq:selected-donor-rational-period}
 D.\mathrm{div}.\mathrm{Bk}\,w(e_{\nu,w}x,e_R y)
     = B_{K,w}(\phi_x)(y).
\end{equation}
In particular the word's norm and the base and twisted periods
are all computed with the same rational comparison.
\releasecite{Development/AbelianVarietyTheory/Proofs/IwasawaDecomposition/ChKerCCongr.lean}{137}{The placed identity for every tuple} and
\releasecite{Development/AbelianVarietyTheory/Proofs/IwasawaDecomposition/RationalPairingEqColumnPairing.lean}{60}{The identity on the descended row}.

\paragraph{The local lattices.}
Write \(N_{(v)},H_{(v)}\) for the localizations inside the
rational class space. For a served place \(v\) let \(c_v\) be
the formal semilinear comparator composed with \(\kappa_w\).
It is class evaluation in the prescribed coefficient frame
at the chosen place embedding. Its image of a basis is
linearly independent over the local coefficient extension,
and the images of integral vectors span the full integral
class model there. In a moving frame it is the base
comparison composed with the inverse tuple transport.
The place is served when the good-prime conditions hold and
the set of sampling vertices is nonempty. Write \(N_{P,v}\)
for the integral test module of that service. The selected
lattice is characterized by
\begin{equation}
\label{eq:selected-donor-localizations}
\begin{aligned}
 v\nmid T,\ {\rm served}: &&N_{(v)}
   &=\{\phi\in V_{\rm ram}^{\rm class}:c_v(\phi)\in N_{P,v}\},\\
 v\mid T: &&N_{(v)}
   &=\{\phi\in V_{\rm ram}^{\rm class}:\kappa_w(\phi)\in N_{\lambda,v}\},\\
 v\nmid T,\ {\rm unserved}: &&N_{(v)}&=H_{(v)}.
\end{aligned}
\end{equation}
Here \(N_{\lambda,v}\) at \(v\mid T\) is the rational preimage
under the stage-zero complex embedding of
\[
 T^{\,m\sum_{s\in\mathcal T_v} b_s}
       \Bigl(\prod_{s\in\mathcal T_v} U_s\Bigr)^m
         \operatorname{emb}_0(S_{\rm ram}),
\]
where \(S_{\rm ram}\) is the formal integral ramified lattice
at \(v\) and \(\mathcal T_v\) consists of the selected top
places if the record marks \(v\) as seen, and is empty
otherwise. The \(U_s\) sum the placed up-neighbours with
coefficient transport. The local up-lattice
\(\Lambda_P+\pi_P^{-1}\mathcal O_P x\) recovers the residual
line of \(x\), and the residual functional condition
\(f_P(x)=1\) fixes its representative. Each term therefore
has the same neighbour, multiplicity and action in the
evaluation on classes. The seen-place label fixes the top
set and the exponents for this word.
\releasecite{Formalism/LocalRingCardinality/DataLatDefs.lean}{13}{The comparator and its integral image},
\releasecite{Formalism/LocalToGlobal/GroupLattices/ComparisonMaps/ComparisonZeroEmbedding.lean}{21}{The moving frame},
\releasecite{Formalism/LocalRingCardinality/DataLatDefs.lean}{42}{The served preimage},
\releasecite{Formalism/LocalRingCardinality/DataLatDefs.lean}{67}{The placed level image},
\releasecite{Formalism/LocalToGlobal/Columns/TiedUpDnA.lean}{42}{The up-lattice},
\releasecite{Development/AbelianVarietyTheory/Proofs/IwasawaDecomposition/LvSameZero.lean}{156}{The decorated up-neighbours} and
\releasecite{Formalism/LocalToGlobal/Columns/AtTOwedS.lean}{266}{The selected top places and exponents}.

Only the places over \(T\) and served primes below \(X'\) can
change the lattice. Each prescribed module is commensurable with
the local \(H_{(v)}\). Choose one integer \(d>0\) with
\(dH_{(v)}\subseteq N_{(v)}\subseteq d^{-1}H_{(v)}\) everywhere.
In \(d^{-1}H/dH\), the local submodules specify one submodule
by its primary components. Its preimage in \(d^{-1}H\) has the
given localizations. Equality at every finite localization gives
global equality. The constructor uses these tests for
\(L.\mathrm{NuLat}\,w\), and the unipotent translate carries this
lattice and its inclusion to the final point. Thus the glued
lattice is precisely \eqref{eq:selected-undivided-donor}.
\releasecite{Development/AbelianVarietyTheory/Proofs/GlobalInterpolation/ExistsLedPart.lean}{119}{Gluing the selected tests},
\releasecite{Development/AbelianVarietyTheory/Proofs/GlobalInterpolation/ExistsLedPart.lean}{187}{The local lengths in the record} and
\releasecite{Development/AbelianVarietyTheory/Definitions/GlobalTwisting/SetUnipotentTranslate.lean}{372}{Carrying the record to the divided point}.

\paragraph{Local length off the level.}
At \(v\mid p\), let \(\alpha_{1v}\le\cdots\le\alpha_{d_r v}\)
be the elementary-divisor exponents of \(N_v\) in \(H_v\) over
the local DVR. The length recorded in the family is normalized
by \(e_v=\operatorname{ord}_v(p)\):
\begin{equation}
\label{eq:selected-donor-raw-length}
 \operatorname{length}_{\mathcal O_v}(H_v/N_v)
     =\sum_j\alpha_{jv}=e_v\ell_{wv}.
\end{equation}
At a served prime off \(T\), write \(M_g\) and \(V_g\) for the
integral class model and rational ramified space over the
extension DVR. For a local lattice \(A\) put
\[
 Y(A)=\sum_B p^{\sum_{s\notin B} S_s}
                  \Bigl(\prod_{s\in B}y_s^m\Bigr)A,\qquad
 S=\max_s S_s,\qquad Y=Y(M_g).
\]
The sum runs over subsets of sampling vertices and uses their
prescribed operators \(y_s\) and integer exponents \(S_s\ge0\).
The tested module is \(N_{P,v}=V_g\cap Y\cap p^{\tau_v}M_g\).
The empty product contains a power of \(p\) times \(M_g\).

The all-but estimate on \(M^{\rm na}\) uses a leading exponent
\(\tau'_v\) and
\(p^{\varepsilon_v}M^{\rm na}\subseteq M_g\subseteq M^{\rm na}\).
Multiplication by \(p^{\varepsilon_v}\) sends
\(Y(M^{\rm na})\) into \(Y(M_g)\), so the new leading exponent is
\(\tau_v=\tau'_v+\varepsilon_v\). For every \(\beta\ge0\) the
image \(Q_\beta\) of \(p^{\tau_v+\beta S}M_g\) modulo \(Y\)
has at most \(d_M\Sigma^{\beta+1}/(\beta+1)!\) generators,
where \(\Sigma=m\sum_s\kappa_s\le1/2\) in the served range.
Indeed, for each set \(B\) with \(\beta+1\) vertices the walk
allows at most \(d_M\prod_{s\in B}\kappa_s\) exceptions at each
of \(m^{\beta+1}\) exponent choices. Summing over \(B\) gives
the factorial bound. The image is a submodule of the quotient
generated by the exceptions, and over a DVR a submodule needs
at most as many generators. Since \(Q_{\beta+1}=p^S Q_\beta\),
the length of \(Q_\beta/Q_{\beta+1}\) is bounded by
\(\operatorname{ord}_g(p)S d_M\Sigma^{\beta+1}/(\beta+1)!\).
For \(S>0\) these images eventually vanish by the empty-product
term. When \(S=0\) that term gives \(Y=M_g\), so the test is
\(p^{\tau_v}(M_g\cap V_g)\). Summing the layers and using
\(e^\Sigma-1\le2\Sigma\) bounds the length of
\(M_g/(Y\cap p^{\tau_v}M_g)\) by
\(\operatorname{ord}_g(p)d_M(\tau_v+\iota_v)\) with
\(\iota_v=2S\Sigma\). The ramified quotient injects into this
quotient.
\releasecite{Formalism/SharpRestriction/GrowthBounds/IsResidueClassZero.lean}{464}{The tested module},
\releasecite{Formalism/SharpRestriction/GrowthBounds/SelectionLemmas.lean}{41}{The all-but length bound} and
\releasecite{Formalism/SubspaceCounting/SubspaceCountingProofs/WalkSharp.lean}{427}{The walk count}.

The comparator gives \(M_g\cap V_g=\mathcal O_g c_v(H_{(v)})\)
by extension of a basis adapted to the saturated inclusion.
Rational spanning for the tested module gives
\(N_{P,v}=\mathcal O_g c_v(N_{(v)})\).
Scalar extension multiplies raw lengths by the relative
ramification index \(e'\), and \(\operatorname{ord}_g(p)=e'e_v\).
The containment in \(p^{\tau_v}M_g\) contracts by the integral
basis comparison to \(p^{\tau_v}H_{(v)}\). It follows that
\begin{equation}
\label{eq:selected-donor-off-level-index}
 \alpha_{1v}\ge e_v\tau_v,\qquad
 \ell_{wv}\le d_M(\tau_v+\iota_v)\qquad(v\nmid T).
\end{equation}
In the concrete local generators one can take
\(M\cap c_v^{-1}(Y)\) as the full-column submodule at this place.
It contains a power of \(p\) times \(M\), and after localization
its intersection with \(p^{\tau_v}M_{(v)}\) and
\(V_{\rm ram}^{\rm class}\) is exactly \(N_{(v)}\) by
\eqref{eq:selected-donor-localizations}. At an unserved place
off \(T\) the length and both parameters vanish.
\releasecite{Formalism/LocalToGlobal/GroupLattices/LinkGConcreteNoKer.lean}{28}{Rational descent on the tied comparison},
\releasecite{Formalism/LocalRingCardinality/Lemmas.lean}{1061}{Scalar extension on the two lattices},
\releasecite{Development/AutomorphicArithmetic/ThetaCorrespondence/ThetaPairLifts/ThirdPlaceSeries/LinkG.lean}{23}{The selected service} and
\releasecite{Formalism/LocalToGlobal/Columns/FiniteSumNonzeroLargeDepth.lean}{1315}{Concrete local generators}.

\paragraph{Local length at the level.}
Write \(\Lambda\) for the stage-zero lattice over \(v\mid T\)
and \(V_{{\rm ram},v}\) for the image of \(H_w\otimes_{\mathcal O_K}K_v\)
under \(\kappa_w\). Let
\(\Lambda_r=\Lambda\cap V_{{\rm ram},v}\) and \(g_s=T^{b_s}U_s\).
The auxiliary lattice \(A_s\) satisfies
\(T^{b_s}\Lambda\subseteq A_s\),
\(U_sA_s\subseteq\Lambda\) and
\(U_s^*\Lambda_r\subseteq A_s\cap V_{{\rm ram},v}\).
Together with \(U_sU_s^*=U_s^*U_s=T^{k-1}\) on the rational
ramified space, these relations give
\[
 g_s\Lambda_r\subseteq U_s(A_s\cap V_{{\rm ram},v})\subseteq\Lambda_r,
 \qquad
 T^{k-1}\Lambda_r\subseteq U_s(A_s\cap V_{{\rm ram},v}).
\]
Split the index at \(U_s(A_s\cap V_{{\rm ram},v})\).
Injectivity of \(U_s\) on the ramified space and the
auxiliary index give
\[
 \operatorname{length}(\Lambda_r/g_s\Lambda_r)
    \le e_v d_M\bigl(k-1+(a_s+kb_s)/k\bigr).
\]
The twist bound is \((a_s+kb_s)/k\le a_{\max}/k+2W_s\),
with \(a_{\max}\le2a_E\), \(W_s\le3b'\) and \(k\le b'\).
Thus \(k-1+(a_s+kb_s)/k\le 2a_E/k+7b'\).
Lengths add for the preserving injective products and their
\(m\)-th powers over at most \(N_T\) top places. Use
\(d_r\ge d_M/2\) beyond the common threshold,
\(\lambda_v=\log(\mathrm N\mathfrak p_v)/[K:\mathbf Q]\), and
\(\sum_{v\mid T}e_v\lambda_v=\log T\). This gives
\begin{equation}
\label{eq:selected-donor-level-index}
 \tau_v=0,\quad \alpha_{1v}\ge0\quad(v\mid T),\qquad
 \frac2{d_r}\sum_{v\mid T} e_v\lambda_v\ell_{wv}
       \le 28N_Tm(a_E/k+b')\log T.
\end{equation}
Here \(\tau_v=0\) by the support of the service. The raw length
in \eqref{eq:selected-donor-raw-length} is the length of the
Hecke image via the same inclusion \(j_w\). The fitted family
bound uses coefficient \(36\), and either bound gives the
stated absolute-constant allowance.
\releasecite{Formalism/LocalToGlobal/GlobalTheory/GL6/CongruenceNoKernelOfHeckeRamified.lean}{44}{The one-operator index},
\releasecite{Formalism/SharpRestriction/GlobalSchurBounds/FiniteModelOfAdmissibleTransports.lean}{1706}{The twist bound},
\releasecite{Formalism/LocalRingCardinality/Lemmas.lean}{1521}{The raw level length},
\releasecite{Formalism/LocalToGlobal/GlobalTheory/PlacesMin.lean}{101}{The support and integrality of \(\tau_v\)} and
\releasecite{Formalism/SharpRestriction/GlobalSchurBounds/LinkMin.lean}{587}{The fitted weighted allowance}.

These localizations and \eqref{eq:selected-donor-placed-period}
apply the valuation test to the very generators used in the
period descent, at the base tuple and at every finite-place tuple.
The unit in the \hyperlink{proof.local-unit-comparison}{moving-tuple comparison}
preserves their valuations. At the level prime, at the two
auxiliary primes and at the remaining primes, the corresponding
integral tests give \eqref{eq:period-gain-valuation}. Hence the
fractional-ideal division there maps \(I_wN_w\) into the dual of
the descended integral row lattice. Its rational pairing is
\eqref{eq:selected-donor-rational-period}, and its source norms
are those of \eqref{eq:selected-donor-isometry}. The local
indices above thus enter the same slopes whose map height is
computed in the \hyperlink{proof.divided-map-heights}{quotient row metric}.
\releasecite{Formalism/LocalToGlobal/GroupLattices/AdmissibleAtGeneratorPlaces.lean}{214}{The three valuation cases} and
\releasecite{Formalism/LocalToGlobal/Divisibility/RecBridge.lean}{631}{The test and division on this source}.

For an independently prescribed lattice \(A\) and a rational
comparison \(\Phi:A_K\to V_{\rm ram}^{\rm class}\), identification
with \(N_w\) requires \(\Phi(A)_v=(N_w)_v\) at every finite
place. In local Smith coordinates this means all the signed
exponents relative to \((N_w)_v\) vanish. The same \(\Phi\)
must carry the Hermitian forms at every embedding and transport
the base and moving periods.
